% SPDX-License-Identifier: Apache-2.0
\section{Type Parameters Given by Name}
\label{sec:x-named-params}

Rust has no named type arguments.
A struct of two type parameters with overlapping bounds, such as an
address and a data word, takes its arguments in either order, and a
swap compiles.
\code{\#[named\_type\_params]} on the struct makes a marker for each
type parameter, named for it, and every use names each argument with
its marker: \code{Bus<Addr<u32>, Word<u64>>} (\#1899).
A marker in the wrong place, or a bare type where a marker belongs, is
refused by the compiler, since the marker of one parameter does not
fill the role of another.
The fields are still the plain types the use named, so \code{show}
takes a \code{u32} and a \code{u64}, and the markers take no room:
the struct is as large as its fields.

Each such struct also has a macro of its own name, which takes every
parameter by name and in any order, \code{Bus![Word = u64, Addr =
u32]}, and refuses a parameter given twice, one the struct does not
have, or one left out that has no default.
That is how a const parameter is given by name: stable Rust cannot read
a number out of a marker, so \code{Lanes![COUNT = 4, WIDTH = 8]} is
the way to write \code{Lanes<8, 4>} with each number named.
The positional form still compiles for a const parameter, so the macro
is a discipline there and not a guarantee.

Two structs that name a parameter alike share its marker:
\code{named\_markers!(pub Addr, pub Word)} declares the markers once,
and \code{Bus} and \code{Beat} each take them with
\code{\#[named\_type\_params(reuse(..))]}, their own bounds stated
for themselves (\#1914).

The runtime's \code{Signal} is written this way, so a wire of eight
bits is \code{Signal<Data<U<8>>>}, and one in a named clock domain
\code{Signal<Data<U<8>>, Domain<C>>}.

\lstinputlisting[caption={\code{ex\_named\_params.rs}. Run with \code{bazel run //lib/examples:ex\_named\_params}.},label={lst:x-named-params}]{lib/examples/ex_named_params.rs}

\lstinputlisting[style=ascii,caption={What \code{ex\_named\_params} printed when the build ran it.},label={lst:x-named-params-out}]{out_named_params.txt}
